\documentclass[landscape,twocolumn,leqno,a4paper,10pt]{report} %leqno affiche les numéros d'équations à gauche

\usepackage[a4paper,top=0.4cm,bottom=1cm,left=0.4cm,right=0.4cm,marginparwidth=0.5cm,footskip=0.6cm]{geometry}



%----------------- Codage de la langue et de la police--------------------------------------------------------------------------
\usepackage[french]{babel} 		% Permet de spécifier au compilateur que l'on désire taper son document en français.
\usepackage[utf8x]{inputenc} 	% Permet de taper directement à l'écran, dans le fichier .tex, les caractères accentués.
\usepackage[T1]{fontenc} 		% Permet juste d'utiliser une nouvelle norme LATEX concernant le codage des caractères.
%\usepackage{lmodern}		% Pack de police --> Pas besoin puisque j'utilise Fourier

%------------------------ Packages utiles-------------------------------------------------------------------------------------------

\usepackage{graphicx}	% Pour insérer des images
\usepackage{caption}		% Légende

\newcommand\hmmax{0}		% Alors c'est une astuce assez mystérieuse pour dépasser la limite de polices latex.
							% En effet, le package bm est lourd et avec tout ce que l'on met, on obtient un message d'erreur : Too many math alphabets used in version normal.			
\newcommand\bmmax{0}		% Cette astuce règle le problème, il faut tout de même mettre ces compteurs avant le package bm

\usepackage{amsmath,amssymb}	% Permettent d'obtenir des caractères mathématiques (de toute beauté !).
\usepackage{amsthm}		% Pour pouvoir faire les newthm
\usepackage{mathtools}	% Corrige des bugs dans amsmath et le complète
\usepackage{numprint}	% Pour les espacements dans les nombres : \numprint{10000} = 10 000
\usepackage{fourier}		% REMARQUE : [upright] impose aux majuscules d'être droite dans le mode mathématique 
\usepackage{pifont} 		% Importe une large gamme de symboles, on utilise \ding{numéro}
\usepackage{xargs} 		% Permet de définir des macros à plusieurs paramètres optionnels.
\usepackage{fancyhdr} 	% Package pour les entêtes et bas de page
\pagestyle{fancy} 			% Accompagne \usepackage{fancyhdr}
\usepackage{bm} 		% Permet de mettre en gras du contenu mathématiques via \bm{..}
\usepackage{mathrsfs} 	% Pour les lettres cursives
\usepackage{enumitem} 	% Pour pouvoir personnaliser les listes

\frenchbsetup{StandardLists=true} % Afin d'éviter tout conflit entre les extensions enumitem et french babel. 

\usepackage{sectsty}	% Pour colorer les titres
\usepackage{titlesec} 	% Pour personnaliser les titres
\usepackage{setspace} 	% Pour changer la taille des interlignes
       	               
\usepackage{array,multirow,makecell}
\usepackage{hlist}		% Package Christian Tellechea sur les listes numérotées gauche vers droite
\setdefaulthlist{left margin = 2mm, label sep= 0.2mm}	% réglages package hlist
				

%------------ Les packages de dessin sous pstricks : il faut changer la compilation de latex pour compiler ces dessins------------------
% Options\Configurer Texmaker\Compil rapide\ Latex+dvips+ps2pdf+Voir PDF

\usepackage[pdf]{pstricks} 	% Pour faire des dessins
\usepackage{pstricks-add,pst-plot,pst-tree,pst-node,pst-math,pst-eucl,pst-func,pst-3dplot}% ,pst-text,pst-eps,pst-fill,pst-grad
\usepackage{multido}		% Pour faire des multido
\usepackage{lastpage} 	% Pour le compteur de la dernière page, on appelle alors celle-ci avec lastage



\usepackage[colorlinks=true, 
			%allcolors=blue
			]{hyperref}		% ATTENTION : doit être importé après la majorité des packages sauf tabularx et breakrules (par exemple)
\hypersetup{			% Paramétrage des hyperliens
	colorlinks = true,   		% Colorise les liens
	breaklinks = true,		% Permet les retours à la ligne pour les liens trop longs
	urlcolor = blue,		% Couleur des hyperliens
	linkcolor = black,		% Couleur des liens internes aux documents (index, figures, tableaux, équations, ...)
	citecolor = blue		% Couleur des liens vers les références bibliographiques
	}




%%%%%%%%%%%%%%%%%%%%%%%%%%%%%%%%%%%%%%%%%%%%%%%%%%
%------------------------- Définition des commandes -----------------------------------------------------------------------------
%%%%%%%%%%%%%%%%%%%%%%%%%%%%%%%%%%%%%%%%%%%%%%%%%%


%--------------------- Raccourcis fontes--------------------------------------------------------------
% REMARQUE : le package fourier met la même police pour mathcal et mathscr, on remet donc mathcal normalement et on utilise la définition de fourier (un peu plus grasse) pour mathscr
\DeclareMathAlphabet{\mathcal}{OMS}{cmsy}{m}{n}	
\DeclareMathAlphabet{\mathscr}{FMS}{futm}{m}{n}

%---------------------------- Raccourcis des espaces--------------------------------------------------------------------------------------------
\newcommand\n{\noindent}
\newcommand\q{\quad}
\newcommand\qq{\quad\quad}

%--------------------- Raccourcis en mode mathématiques-------------------------------------------------------------------------------------
\newcommand{\et}{\text{ et }}
\newcommand{\ou}{\text{ ou }}
\newcommand{\SSI}{\Longleftrightarrow}
\let\si\relax
\newcommand{\si}{\text{ si }}
\newcommand{\sinon}{\text{ sinon }}
\renewcommand{\epsilon}{\varepsilon}		% epsilon beaucoup plus joli
\renewcommand{\geq}{\geqslant} 
\renewcommand{\leq}{\leqslant}
\newcommand{\e}{\mathrm{e}}
\renewcommand{\textbf}[1]{\begingroup\bfseries\mathversion{bold}#1\endgroup}		% Très important, permet de passer en gras même avec le mode mathématiques

\newcommand{\Mod}[1]{\allowbreak\if@display\mkern5mu\else\mkern5mu\fi\left[#1\right]}	% Pour modulo entre crochets [.] avec un espacement convenable 


%---------------------------------- Raccourcis calcul diff----------------------------------------------------------------------------------------
\renewcommand{\d}{\;\mathrm{d}}	% Pour le d droit de l'intégrale car ce n'est pas une variable...


%-------------------------------- Raccourcis des ensembles----------------------------------------------------------------------------
\newcommand\R{\mathbb{R}}
\newcommand\C{\mathbb{C}}
\newcommand\K{\mathbb{K}}
\newcommand\Q{\mathbb{Q}}
\newcommand\D{\mathbb{D}}
\newcommand\Z{\mathbb{Z}}
\newcommand\N{\mathbb{N}}
\newcommand\U{\mathbb{U}}


%------------------------------- Dimensions intelignes, items...-----------------------------------------------------------------------------------------------
\onehalfspacing	% Pour que l'interligne soit de 1.5. Pour que l'interligne soit de la valeur que l'on souhaite, de base : \setstretch{1}

\setlist[1]{font=\bfseries,itemsep=0.2mm,leftmargin=6mm}
\setlist[2]{font=\bfseries,itemsep=0.2mm,leftmargin=4mm}
\setenumerate[1]{font=\bfseries,itemsep=0.2mm,leftmargin=6mm}
\setenumerate[2]{font=\bfseries, label=\alph*., leftmargin=4mm}
%%\setenumerate[profondeur]{font=\bfseries, label=\ding{221} \color{blue}} permet de changer en a) b) les sous-items d'une liste enumerate
%% itemsep = correspond à l'interligne entre chaque item
% de base : labelsep=2mm et permet de changer la distance entre l’étiquette et le texte des éléments
% leftmargin =4mm spécifie la marge à gauche de la liste
\AtBeginDocument{\renewcommand{\labelitemi}{\textbullet}}	% Pour que les puces par défaut soit des points.
%\AtBeginDocument{\renewcommand{\labelitemi}{\ding{86}}}	% Pour que les puces par défaut soit des étoiles.
\AtBeginDocument{\renewcommand{\labelitemii}{$\ast$}}		% Pour que les puces par défaut pour une sous-liste soit des astérisques.
\newcommand{\itemmm}{\item[$\ast$]}				% \itemmm devient l'item \ast
\newcommand{\itemmmm}{\item[\footnotesize$\circledcirc$\normalsize]}
\newcommand{\itemmmmm}{\item[\footnotesize$\astrosun$\normalsize]}
\newcommand{\itemmmmmm}{\item[\footnotesize$\multimap$\normalsize]}
\newcommand{\itemmmmmmmmm}{\item[$\divideontimes$]}
\newcommand{\etoile}{\item[$\star$]}
\newcommand{\etoilee}{\item[$\bigstar$]}
\newcommand{\losange}{\item[$\diamondsuit$]}
\newcommand{\carre}{\item[\LARGE$\centerdot$]}
\newcommand{\carreombre}{\item[\ding{112}]}

\setlength{\columnsep}{30pt} 		% Permet d'écarter les deux colonnes du twocolumn de 30pt, pratique pour écarter le texte de la ligne centrale
\setlength{\columnseprule}{0.5pt}	% Affiche une ligne verticale entre chaque item et entre chaque page du twocolumn
						% \columnseprule=0cm 	% Permet d'enlever le trait vertical entre chaque colonne
%\renewcommand\tabenumsep{0.5cm\hskiplem}	% Ecarte de 0.5cm les items de tabenum, sinon ils sont trop collés...


\newlength{\tagmarginsep} 		% On crée une longueur
\setlength{\tagmarginsep}{0cm}		% Permet de régler la position d'un tag pour une équation mettre avant l'équation : \setlength{\tagmarginsep}{x cm} pour que le tag se rapproche de l'équation
\everydisplay{\displayindent=\tagmarginsep \displaywidth=\dimexpr\linewidth-2\tagmarginsep}	% Va avec ce qu'il y a juste avant


%-------------------- Définition des couleurs-----------------------------------------------------------------------------------------

\definecolor{violetfonce}{RGB}{126,62,142}
\definecolor{magentafonce}{RGB}{128, 0, 128}
\newrgbcolor{magentafonce}{.5 0 .5}			% Une autre façon de définir une couleur, permet surtout d'utiliser \magenta
\definecolor{violet1}{cmyk}{0.5,1,0,0}
\definecolor{violet2}{cmyk}{0.24998, 1, 0, 0}
\definecolor{violet3}{cmyk}{0, 1, 0.25, 0}



\newcommand\tred{\textcolor{red}}
\newcommand\tredf{\textcolor{rougefonce}}
\newcommand\tdred{\textcolor{deepred}}
\newcommand\tblue{\textcolor{blue}}
\newcommand\tbluef{\textcolor{bleufonce}}
\newcommand\tgreen{\textcolor{vertfonce}}
\newcommand\torange{\textcolor{orange}}
\newcommand\tmagenta{\textcolor{magentafonce}}


%------------------ Changement paramètres des parties, sous parties---------------------------------------------------------


\sectionfont{\color{violet1}}			% Utilisation du package sectsty
\subsectionfont{\color{violet2}}
\subsubsectionfont{\color{violet3}}


%--------------------------------- Raccourcis probabilités--------------------------------------------------------------------------------
\newcommand{\Prob}{\mathbb{P}}
\newcommand{\E}{\mathbb{E}}
\newcommand{\normale}[2]{\mathcal{N}\left(#1;#2 \right)}	% Pour la loi Normale : 	$\normale{\mu}{\sigma^2}$
\newcommand{\Unif}[1]{\mathcal{U}\left(#1\right)}			% Pour la loi uniforme :	$\Unif{[a;b]}$
\newcommand{\Ber}[1]{\mathcal{B}\left(#1\right)}			% Pour la loi de Bernoulli : $\Ber{p}$
\newcommand{\Bino}[2]{\mathcal{B}\left(#1;#2\right)}		% Pour la loi Binomiale :	$\Bino{n}{p}$
\newcommand*{\binomial}[2]{% 
  \ensuremath{ 
    \begin{pmatrix} 
      #2\\ 
      #1
    \end{pmatrix}}}				% Pour le coefficient binomial : $\binom{k}{n}$
\renewcommand{\binom}{\binomial} 	% La commande \binom existait déjà mais n'était pas pratique



%----------------------------------- Opérateurs---------------------------------------------------------------------------------------------------

\newcommand{\barre}[1]{\overline{#1}}  									% Barre au dessus du texte, pratique pour le conjugué des complexes
\newcommandx{\Sum}[3][3=\white]{{\displaystyle\sum\limits_{\substack{#1 #3}}^{#2}}}		% $\Sum{k=0}{n}[<option>]$ où [<option>] est placée à la suite de k=0 si l'on tape de façon classique, et en-dessous si l'on met [\\ k\neq k_0]
															% On peut juste taper : $\Sum{k=0}{n}$ et cela fonctionne
															% \displaystyle impose de gérer les tailles comme en mode mathématique : ici on aura donc toujours de grande somme que ce soit en math comme dans le texte
\newcommandx{\Prod}[3][3=\white]{{\displaystyle\prod\limits_{\substack{#1 #3}}^{#2}}}		% $\Prod{k=0}{n}[<option>]$	idem pour l'option
\newcommand{\Int}[4]{{\displaystyle\int_{#1}^{#2} #3  \d #4}}   					% $\Int{0}{\pi}{\cos(t)}{t}$
\newcommandx{\Lim}[3][3=\white]{\lim\limits_{\substack{#1 \to #2 #3}}}		% $\Lim{x}{0}[<option>]$	idem pour l'option


%------------------------------------ Délimiteurs--------------------------------------------------------------------------------------------------
%\newcommand{\abs}[1]{\left\lvert #1 \right\rvert}			% Valeur absolue | |
%\newcommand{\norm}[1]{\left\lVert #1 \right\rVert}			% Doubles barres verticales : norme d'un vecteur || ||

%--------------------------------------- Vecteurs----------------------------------------------------------------------------------------------------
 
\newcommand{\Oij}{\left(O;\vec{\imath},\vec{\jmath}\right)}
\newcommand{\Oijk}{\left(O;\vec{\imath},\vec{\jmath},\vec{k}\right)}
\newcommand{\Ouv}{\left(O;\vec{u},\vec{v}\right)}
\newcommand{\Ouvw}{\left(O;\vec{u},\vec{v},\vec{w}\right)}

%------------------------------------ Systèmes d'équations-------------------------------------------------------------------------------------------
% Remarque \syst et \syste ne font pas d'alignement mais ils sont bien, si l'on veut de l'alignement, on utilisera \systeme du package "systeme"
\newcommand*{\syst}[2]{%			% Systèmes 2 éq. à ? inconnues : 	\syst{ax+by+c=0}{a'x+b'y+c'=0}
  \left\{\begin{array}{ll} 
        #1\\ 
        #2 
      \end{array}  
    \right. }
 
\newcommand*{\syste}[3]{% 			% Systèmes 3 éq. à ? inconnues : 	\syste{ax+by+c=0}{a'x+b'y+c'=0}{a''x+b''y+c''=0}
  \left\{\begin{array}{ll} 
        #1\\ 
        #2\\
        #3 
      \end{array}  
    \right.}



%------------------------------------------ Autre----------------------------------------------------------------------------------------------------------
\newcommand{\boite}[2]{\noindent\parbox[]{#1 cm}{#2} }  									% Boite autre façon de faire minipage en gros


%---------------------------- En-têtes et bas de pages----------------------------------------------------------------------------
% A chaque fois : *head* pour le haut et *foot* pour le bas
% L = left, C = center, R = right 
% L = toutes les pages, LE = toutes les pages paires, LO = toutes les pages impaires, idem CE, CO, ...


\newcommand{\fcours}[1]{								% f pour fancy				   
	\pagestyle{fancy}									% Efface le style qui a pu être mis avant
	\renewcommand{\headrulewidth}{0pt} 
	\renewcommand{\footrulewidth}{0.5pt} 
	%\fancyhead[L]{}									
	%\fancyhead[C]{}
	%\fancyhead[R]{}									
	\fancyfoot[L]{\vspace{-1mm}\footnotesize \textit{#1}} 			% Classe : <Seconde> 
	\fancyfoot[C]{\vspace{-1mm}\footnotesize \thepage/\pageref{LastPage}}
	\fancyfoot[R]{\vspace{-1.8mm}\textit{\footnotesize C. Marguerite}}	% N'apparaîtra pas si on met les marges avec moins de bottom=1.3cm dans le package geometry
 }												% Fin de la commande
					%% Mettre au début du document : \fcours{Seconde} par exemple
					%% \usepackage[a4paper,top=0.4cm,bottom=1cm,left=0.4cm,right=0.4cm,marginparwidth=0.5cm,footskip=0.6cm]{geometry} 




%---------------------------- Réglages chapitres + table des matières------------------------------------------------------
\titleformat{\part}[display]
{\centering\normalfont\Huge\bfseries}
{\partname}
{20pt}
{\rule{6cm}{1pt}\raisebox{-2.1pt}[10pt][10pt]{\quad\decofourleft\decoone\decofourright\quad}\rule{6cm}{1pt} ~\\\vspace{20pt}}

\makeatletter
\renewcommand\chapter{\global\@topnum\z@
                    \@afterindentfalse
                    \secdef\@chapter\@schapter}
\def\@topnewpage[#1]{#1}
\makeatother										% Pour avoir le chapitre sur une demie page

\titleformat{\chapter}{\centering\normalfont\Large\bfseries\scshape\color{magentafonce}}{}{1em}{}
\titlespacing*{\chapter}{0pt}{-18pt plus 0pt minus 2pt}{0pt plus 2pt minus 2pt}		% Réglages espacements (à gauche, au-dessus, en-dessous) de la commande \chapter

\makeatletter
\newcommand{\twocolumntoc}{		
  \section*{\huge\centering\contentsname
    \@mkboth{%
      \MakeUppercase\contentsname}{\MakeUppercase\contentsname}}%
  \@starttoc{toc}%
}
\makeatother										% Permet de mettre la table des matières en twocolumn 
												% Mettre après \begin{doc..} \twocolumntoc \clearpage 		% Remplace \tableofcontents par \twocolumntoc en twocolumn

\makeatletter
\renewcommand{\thesection}{\@arabic\c@section}
\makeatother										% Enlève la numérotation des chapitres devant les sections

\setcounter{tocdepth}{3} 								% Permet d'afficher dans la table des matières suivant une certaine profondeur : -1=part, 0=chapter, 1=section, ...
\setcounter{secnumdepth}{3}							% Pour avoir la numérotation des subsubsections avec les classes book et report

\newcommand{\nocontentsline}[3]{}
\let\oldaddcontentsline\addcontentsline
\newcommand{\tocless}[2]{\let\addcontentsline\nocontentsline #1{#2}\let\addcontentsline\oldaddcontentsline}	% Pour retirer une section de la table des matières : \tocless\section{...}





\setstretch{1}
\begin{document}
\fcours{Concours général}


\chapter{Concours général des lycées - Session 2003}


\n \textit{Le problème étudie des configurations du plan rapporté à un repère orthonormé $(O, \vec{u}, \vec{v})$. On pourra aussi se placer dans le plan complexe associé, i étant l'affixe du point de coordonnées $(0,1)$}.\\

\n \textit{On appelle \textbf{triangle} tout ensemble de \textbf{trois points non alignés} du plan.}

\section{Questions préliminaires}

\begin{enumerate}
	\item Soit $A B C$ un triangle et $M$ un point quelconque du plan. Montrer que :
	$$\overrightarrow{M A} \cdot \overrightarrow{B C}+\overrightarrow{M B} \cdot \overrightarrow{C A}+\overrightarrow{M C} \cdot \overrightarrow{A B}=0$$
	En déduire que les trois hauteurs du triangle $A B C$ sont concourantes en un point $H$ appelé \textbf{orthocentre} de ce triangle.
	\item Soit $A B C$ un triangle, $\Omega$ le centre de son cercle circonscrit et $H$ le point tel que $\Omega H=\Omega A+\Omega B+\Omega C$.\\
	Démontrer que $H$ est l'orthocentre du triangle $A B C$.
\end{enumerate}

\n \textit{Étant donnée une partie $X$ du plan, supposée non incluse dans une droite, on note $\mathcal{H}(X)$ l'ensemble des orthocentres des triangles dont les sommets appartiennent à $X$.\\
On dira qu'une partie $X$ du plan est \textbf{orthocentrique} si elle n'est pas incluse dans une droite et si $\mathcal{H}(X)$ est inclus dans $X$, c'est-à-dire si tout orthocentre d'un triangle de points de $X$ appartient à $X$.}

\section{Première partie}

\begin{enumerate}
	\item Déterminer les parties orthocentriques à 3 éléments.
	\item Déterminer les parties orthocentriques à 4 éléments.
	\item Soit $X$ un ensemble de quatre points d'un cercle et $Y=\mathcal{H}(X)$.
	\begin{enumerate}
		\item Montrer que $Y$ se déduit de $X$ par une transformation simple.
		\item Déterminer $\mathcal{H}(Y)$.
	\end{enumerate}
	\item \begin{enumerate}
		\item Soit $\Gamma$ un cercle de rayon strictement positif; déterminer $\mathcal{H}(\Gamma)$.
		\item Soit $D$ un disque de rayon strictement positif; déterminer $\mathcal{H}(D)$.
	\end{enumerate}
\end{enumerate}

\newpage

\section{Deuxième partie}

\n Dans cette partie, $R$ est un nombre réel strictement positif, $n$ est un entier au moins égal à 2 et $X$ est l'ensemble des $2 n$ sommets d'un polygone régulier inscrit dans le cercle de centre $O$ et de rayon $R$.\\
On considère l'ensemble $\mathcal{T}$ des triangles dont les sommets appartiennent à $X$. On choisit au hasard, avec équiprobabilité, un élément de $\mathcal{T}$.

\begin{enumerate}
	\item Quelle est la probabilité de choisir un triangle rectangle ?
	\item Quelle est la probabilité de choisir un triangle dont les trois angles sont aigus ?
	\item On note $L$ la variable aléatoire qui à tout élément de $\mathcal{T}$ associe le carré de la distance de $O$ à son orthocentre. Déterminer, en fonction de $n$ et $R$, l'espérance de la variable aléatoire $L$.
\end{enumerate}


\section{Troisième partie}


\begin{enumerate}
	\item Soient $a$, $b$, $c$ trois réels tels que $a(b-c) \neq 0$ et $A$, $B$, $C$ les points de coordonnées respectives $(0, a)$, $(b, 0)$, $(c, 0)$.\\
	Calculer les coordonnées de l'orthocentre $D$ du triangle $A B C$.
	\item Soit $X$ la partie obtenue en prenant la réunion d'une droite $\Delta$ et d'un point $M$ n'appartenant pas à $\Delta$.\\
	Déterminer $\mathcal{H}(X)$. Montrer que $\mathcal{H}(X) \cup X$ est une partie orthocentrique.
	\item Soit $X$ une partie orthocentrique incluse dans la réunion des axes $(O, \vec{u})$ et $(O, \vec{v})$ et contenant au moins trois points de $(O, \vec{u})$ distincts de $O$.
	\begin{enumerate}
		\item Montrer que $X$ contient au moins trois points de $(O, \vec{u})$ d'abscisses non nulles et de même signe
		\item Montrer que $X$ contient au moins trois points de $(O, \vec{u})$ d'abscisses strictement positives.
	\end{enumerate}
	\item \begin{enumerate}
		\item Déterminer les parties orthocentriques finies, contenant au plus cinq points et incluses dans la réunion des axes $(O, \vec{u})$ et $(O, \vec{v})$.
		\item Soit $X$ une partie orthocentrique incluse dans la réunion des axes $(O, \vec{u})$ et $(O, \vec{v})$ et contenant au moins six points. Montrer qu'il existe deux suites $\left(x_{n}\right)$ et $\left(x_{n}^{\prime}\right)$ de réels non nuls telles que, pour tout entier $n$, les points de coordonnées $\left(x_{n}, 0\right)$ et $\left(x_{n}^{\prime}, 0\right)$ appartiennent á $X$, et telles que :
		$$\Lim{n}{+\infty} x_{n}=+\infty, \quad \Lim{n}{+\infty} x_{n}^{\prime}=0$$
		Une partie orthocentrique incluse dans la réunion des axes $(O, \vec{u})$ et $(O, \vec{v})$ et contenant au moins six points peut-elle être finie ?
	\end{enumerate}
\end{enumerate}

\newpage

\section{Quatrième partie}

\n \textit{L'objectif de cette partie est la construction de parties orthocentriques remarquables.}

\begin{enumerate}
	\item Soit $k$ un réel non nul et soit $Y$ l'hyperbole d'équation $x y=k$.
	\begin{enumerate}
		\item Soient $A$, $B$, $C$, $D$ quatre points distincts de $Y$, d'abscisses respectives $a$, $b$, $c$, $d$. Montrer que $\overrightarrow{A B}$ et $\overrightarrow{C D}$ sont orthogonaux si et seulement si $a b c d=-k^{2}$.
		\item Soient $A$, $B$, $C$ trois points distincts de $Y$, d'abscisses respectives $a$, $b$, $c$. Déterminer l'orthocentre de $A B C$.
		\item Montrer que $Y$ est orthocentrique.
	\end{enumerate}
\end{enumerate}

\n \textit{Dans toute la suite de la quatrième partie, on considère un entier relatif non nul $q$ et on note $X$ l'ensemble d'équation $x^{2}+q x y-y^{2}=1$.}

\begin{enumerate}\setcounter{enumi}{1}
	\item \begin{enumerate}
		\item Montrer que l'équation $t^{2}-q t-1$ possède deux racines réelles distinctes. Montrer que ces racines sont irrationnelles.
	\end{enumerate}
\end{enumerate}

\n \textit{Dans toute la suite de la quatrième partie, on note $r$ et $r^{\prime}$ ces deux racines et $s$ la similitude définie par la représentation complexe $z \mapsto(1-r i) z$.}

\begin{enumerate}\setcounter{enumi}{2}
	\item[] \begin{enumerate}\setcounter{enumi}{1}
		\item Montrer que $s(X)$ est une hyperbole, d'équation $x y=k$, où $k$ est un réel à déterminer. En déduire que $X$ est un ensemble orthocentrique.
	\end{enumerate}
	\item Soit $G$ l'ensemble des points de $X$ à coordonnées entières et $\Gamma$ l'ensemble des abscisses des éléments de $s(G)$.
	\begin{enumerate}
		\item Vérifier que $\Gamma$ est l'ensemble des nombres réels de la forme $x+r y$, où $x$ et $y$ sont deux entiers tels que $(x+r y)\left(x+r^{\prime} y\right)=1$.
		\item Montrer que $-1 \in \Gamma$; montrer que $r^{2} \in \Gamma$.
		\item Montrer que le produit de deux éléments de $\Gamma$ est élément de $\Gamma$ et que l'inverse d'un élément de $\Gamma$ est élément de $\Gamma$. Montrer que $\Gamma$ possède une infinité d'éléments.
	\end{enumerate}
	\item Déduire de ce qui précède que l'ensemble $G$ des points à coordonnées entières de $X$ est une partie orthocentrique infinie.
\end{enumerate}


\section{Cinquième partie}

\n On note $Y_{1}$ l'hyperbole d'équation $x y=1$ et $Y_{0}$ l'ensemble d'équation $x y=0$, c'est-à-dire la réunion des axes $(O, \vec{u})$ et $(O, \vec{v})$.\\
On admet le résultat suivant :\\
<< \textit{Étant donnés quatre points $A$, $B$, $C$ et $D$ du plan, il existe une similitude $s$ telle que $s(A)$, $s(B)$, $s(C)$ et $s(D)$ appartiennent tous à $Y_{1}$ ou bien appartiennent tous à $Y_{0}$} >>.\\

Soient $A_{0}$, $B_{0}$, $C_{0}$ et $D_{0}$ quatre points, trois à trois non alignés, et soit $X_{0}=\left\{A_{0}, B_{0}, C_{0}, D_{0}\right\}$.\\
On définit par récurrence $X_{n+1}=\mathcal{H}\left(X_{n}\right)$ pour tout entier naturel $n$. On suppose qu'il existe un entier $n$ strictement positif tel que $X_{n}=X_{0}$ et on note $m$ le plus petit entier ayant cette propriété.

\begin{enumerate}
	\item Montrer que $m=1$ ou $m=2$.
	\item Déterminer les ensembles $X_{0}$ tels que $m=1$, puis ceux tels que $m=2$.
\end{enumerate}


\end{document}